\chapter{}
\label{chap:3}

\section{Умова}

Нехай поліном $f$ утворює базис Грьобнера ідеал $I = \langle f \rangle$ відносно мономіального впорядкування
$\preceq$. Довести, що $\text{LM}_{\preceq}(f)$ є найменшим відносно $\preceq$ елементом множини
$\left\{\text{LM}_{\preceq}(g) \, \vert \, g \in I \setminus {0}\right\}$. Чи є правильним обернене твердження?
(Якщо нi - побудувати контрприклад)

\section{Розв'язання}

\begin{proof}
    ~\par \noindent Оскільки $f$ належить базису Грьобнера ідеалу $I$, то за означенням маємо:
    \begin{equation*}
        \bigl\langle \LM(f) \bigr\rangle = \bigl\langle \LM(g) \, \vert \, g \in I \setminus \{0\} \bigr\rangle.
    \end{equation*}
    Тобто для будь-якого $g \in I \setminus \{0\}$ виконується наступне:
    \begin{equation}
        \label{eq:div}
        \LM(f) \, \vert \, \LM(g)
    \end{equation}
    Тобто $\LM(g)$ -- найстарші мономи, які уторюють ідеал. А також будь-якого допустимого мономіального впорядкування 
    $\preceq$ подільність мономів випливає наступне:
    \begin{equation*}
        m_1 \mid m_2 \, \Longrightarrow \, m_1 \preceq m_2
    \end{equation*}
    Дійсно, якщо $m_2 = m_1 \cdot m'$ для деякого монома $m' \neq 0$, то за властивістю впорядкування $1 \preceq m'$, 
    і тому
    \begin{equation*}
        m_1 = m_1 \cdot 1 \preceq m_1 \cdot m' = m_2
    \end{equation*}
    Ну і тоді з~\eqref{eq:div} випливає, що $\LM(f) \preceq \LM(g)$ для всіх $g \in I \setminus \{0\}$.

    А оскільки $f \in I \setminus \{0\}$ (бо $f \neq 0$, бо він є елементом базису Грьобнера), ми маємо
    $\LM(f) \in M$. Отже, $\LM(f)$ є найменшим елементом множини $M$ відносно впорядкування $\preceq$.
\end{proof}

\noindent Зворотнє твердження звучатиме так:

\textit{<<Якщо $\LM_{\preceq}(f)$ є найменшим елементом множини $M$, то $\{f\}$ є базисом Грьобнера ідеалу $I = \langle f \rangle$.>>}

Мінімальність $\LM(f)$ за порядком $\preceq$ означає лише $\LM(f) \preceq \LM(g)$ для всіх $g \in I \setminus \{0\}$. Однак 
для того, щоб $\{f\}$ був базисом Грьобнера, необхідно щоб виконувалася ознака \emph{подільності}: $\LM(f) \mid \LM(g)$ 
для всіх таких $g$. Ці дві умови не є еквівалентні: з $m_1 \preceq m_2$ не випливає $m_1 \mid m_2$. Тому \textbf{твердження є хибним.}

Напрошується контрприклад (фактично, можна зробити відсилку на пункт Б задачі 1 з третьої практики і на цьому 
завершити, бо там як раз це і проявилося), але, для формальності, я продублюю власне наші думки :)

Розглянемо векторний простір $B_3$, і drl впорядкування $\preceq_{\mathrm{drl}}$ ($x_1 \succ x_2 \succ x_3$) та 
власне поліном
\begin{equation*}
    f = x_1 \oplus x_2 x_3.
\end{equation*}

\noindent Впорядкування мономів у $B_3$ за $\preceq_{\mathrm{drl}}$ має наступний вигляд:
\begin{equation*}
    1 \prec x_3 \prec x_2 \prec x_1 \prec x_2 x_3 \prec x_1 x_3 \prec x_1 x_2 \prec x_1 x_2 x_3.
\end{equation*}

Розглянемо наш поліном $f$. Оскільки $\deg(x_2 x_3) = 2 > 1 = \deg(x_1)$, то маємо, що старший моном $f$ відносно 
drl це $\LM_{\mathrm{drl}}(f) = x_2 x_3$. Припустимо що $\{f\}$ -- базис Грьобнера, тоді обчислимо $(1 \oplus x_2) \cdot f$:
\begin{equation*}
    (1 \oplus x_2) f = (x_1 \oplus x_2 x_3) \oplus x_2(x_1 \oplus x_2 x_3)
    = (x_1 \oplus x_2 x_3) \oplus (x_1 x_2 \oplus x_2 x_3)
    = x_1 \oplus x_1 x_2 = g.
\end{equation*}
Оскільки ми множимо $f$ на інший елемент кільця, то результат $g = x_1 \oplus x_1 x_2 \in I$ (зі старшим мономом $\LM_{\mathrm{drl}}(g) = x_1 x_2$) 
гарантовано належитиме ідеалу $I$.

Але оскільки $x_3 \nmid x_1 x_2$, маємо $x_2 x_3 \nmid x_1 x_2$ -- новий моном не ділить наш базисний, тобто $\LM(f) \nmid \LM(g)$. 
І отримуємо пряме порушення означення базису Грьобнера. Отже, $\{f\}$ не є базисом Грьобнера.

А з іншого боку $\LM(f) = x_2 x_3$ є найменшим елементом $M$, зараз доведемо що це так. Покажемо, що в $I$ немає 
ненульових елементів зі старшим мономом степеня $\leq 1$.

Нехай для деякого $h \in B_3$, $h \neq 0$ визначимо $g = h \cdot f = h(x_1 \oplus x_2 x_3) = h x_1 \oplus h x_2 x_3$. 
Зауважимо, що будь-який моном добутку $h \cdot x_2 x_3$ матиме степінь $\geq 2$ (бо містить принаймні $x_2$ та $x_3$). 
Щоб такий моном скоротився, необхідно, щоб він з'являвся у $h \cdot x_1$. Але всі мономи $h \cdot x_1$ містять змінну 
$x_1$, тоді як моном $x_2 x_3$ (що виникає з $h$, якщо $h = \ldots \oplus 1$) не міститиме $x_1$ і тому він не 
скоротиться з жодним мономом $h \cdot x_1$.

Отже, якщо $1$ є мономом в $h$ (тобто вільний член $h$ ненульовий), то $x_2 x_3$ буде обов'язково присутнім в $g$, і 
тоді $\LM_{\mathrm{drl}}(g)$ має степінь $\geq 2$. А якщо ж $h$ не містить вільного члена, тобто $h(0,0,0) = 0$, то 
кожен моном $h$ містить хоча б одну змінну, і мономи $h x_{2} x_{3}$ матимуть степінь $\geq 3$. Мономи $hx_1$ мають
степінь $\geq 2$. Тому $\LM_{\mathrm{drl}}(g)$ має степінь $\geq 2$.

В обох випадках степінь найстаршого монома $\deg(\LM_{\mathrm{drl}}(g)) \geq 2$. Серед мономів степеня $2$ при drl
виконується наступне: $x_2 x_3 \prec x_1 x_3 \prec x_1 x_2$. А оскільки $\LM(f) = x_2 x_3$ це найменший моном степеня 
$2$, він є найменшим елементом $M$.

\noindent \textbf{Висновок.} Обернене твердження хибне. І контрприклад з $f = x_1 \oplus x_2 x_3 \in B_3$, 
$\preceq\; = \;\preceq_{\mathrm{drl}}$ слугує тому доведенням. Тут $\LM(f) = x_2 x_3$ є найменшим у $M$, але 
$\{f\}$ не є базисом Грьобнера, оскільки існує таке $g = x_1 \oplus x_1 x_2 \in I$, що $\LM(g) = x_1 x_2$, і при 
цьому $x_2 x_3 \nmid x_1 x_2$.
